\begin{table}[ht]
\centering
\caption{Sample construction and empirical units.}
\label{tab:sample_construction_v9}
\begin{threeparttable}
\small
\begin{tabular}{p{.28\linewidth}p{.24\linewidth}p{.38\linewidth}}
\toprule
Object or sample & Size & Unit and use \\
\midrule
Classifier universe & \BPregexNClassified{} & Purchase orders/tender notices classified into procurement regimes. \\
Full BEC pharmaceutical file & \BPnPOIfull{} & Purchase-offer-item observations after linking BEC item records to regime classifications. \\
Analysis sample & \BPnAnalysisSample{} & POI observations satisfying core item and variable restrictions. \\
Winners-only price sample & \BPnegNobs{} & Accepted winning bids used in negotiated-price regressions. \\
Urgent panel & \BPnUTG{} & Administrative and litigated urgent winning bids used for the under-the-gun comparison. \\
Both-regime urgent cells & \BPbothCellN{} & Item-by-year-month cells containing both administrative and litigated urgent purchases. \\
Singleton urgent cells & \BPbothCellNSingleton{} & Comparable urgent cells containing only one urgent regime. \\
Firm-buyer-item triples & \BPutgTripleCountUTG{} triples; \BPutgTripleNUTG{} observations & Same supplier, buyer, and item comparisons used to isolate within-firm pricing. \\
\bottomrule
\end{tabular}
\begin{tablenotes}[flushleft]\footnotesize
\item \textit{Notes:} This table separates the classifier unit from the empirical regression units. POI denotes purchase-offer-item. Price regressions use accepted winning bids, while classifier validation operates at the purchase-order/tender-notice level.
\end{tablenotes}
\end{threeparttable}
\end{table}
